\uuid{QO24}
\titre{Tangente}
\auteur{bodin ; hanani ; mzari}
\organisation{exo7}
\licence{CC BY-NC-SA 4.0}
\datecreate{2020-01-20}
\niveau{L1}
\module{Analyse}
\chapitre{Dérivabilité des fonctions réelles}
\sousChapitre{Autre}
\difficulte{1}
\theme{Dérivées, tangente}
\exo7qcmid{lille-1.derivees.001}

\contenu{
\begin{question}
\qclassification{124.00}
\qtags{motcle=tangente}
Soit $\displaystyle f(x)=\frac{2}{x}$ et $g(x)=2\sqrt{x}$. On note $\mathcal{C}_f$ (resp. $\mathcal{C}_g$) la courbe représentative de $f$ (resp. $g$). Quelles sont les bonnes réponses ?
\begin{answers}
    \good{Une équation de la tangente à $\mathcal{C}_f$ au point $(1,2)$ est $y=-2x+4$.}
    \bad{Une équation de la tangente à $\mathcal{C}_f$ au point $(1,2)$ est $y=-2x+2$.}
    \bad{Une équation de la tangente à $\mathcal{C}_g$ au point $(1,2)$ est $y=x+2$.}
    \good{Une équation de la tangente à $\mathcal{C}_g$ au point $(1,2)$ est $y=x+1$.}
\end{answers}
\begin{explanations}
Une équation de la tangente à $\mathcal{C}_f$ au point $(a,f(a))$ est :
$$y=f'(a)(x-a)+f(a).$$
Ici, $\displaystyle f'(x)=-\frac{2}{x^2}$ et $\displaystyle g'(x)=\frac{1}{\sqrt{x}}$.
\end{explanations}
\end{question}
}
